\chapter{WiFi 驱动：XR829 / SDIO 枚举与固件下载}
\label{chap:wifi}

本章对应当前 miniOS 的实现，说明 100ASK T113-Pro 板上 XR829 WiFi/BT 模组驱动的分层结构、SDIO 枚举与 CMD53 数据通路、SMHC 主机控制器、HIF 寄存器/间接寻址、固件下载全流程，以及固件镜像的打包与加载。

当前已打通的链路：

\begin{itemize}
  \item SDIO 枚举（CMD0/5/3/7 + CMD52 读 CCCR/FBR/CIS）；
  \item CMD53 双向数据通路（PIO，非 DMA）与 HIF 寄存器自检；
  \item 总线提速（400 kHz $\rightarrow$ 12 MHz 默认速度、1-bit）并在高速下复测；
  \item HIF 间接访问（SRAM/AHB 端口）与 DPLL/唤醒序列；
  \item boot\_xr829 载入 + fw\_xr829 通过 32 KiB 下载 FIFO 传输；
  \item 切到 queue（消息）模式，解析 WSM startup indication；
  \item WSM 命令/confirm 串行化、RX indication 分派、TX confirm、MIB 读写；
  \item WSM CONFIGURATION、2.4 GHz 定向主动扫描，以及 beacon/probe response 解析；
  \item WPA2-PSK 所需 PBKDF2/PTK/MIC/AES unwrap 通用密码学模块。
\end{itemize}

主机侧没有接 SDIO 中断（IRQ 73 已声明但全部轮询），整条 WiFi 链路是轮询驱动的。

\textbf{更新}：STA JOIN、WPA2 四次握手（EAPOL）和 IPv4/UDP/DHCP/tftp 都已实现，网络栈和驱动已按 Linux 分层重新组织（cfg80211 / mac80211 / supplicant / xradio 驱动拆分），网络栈细节见第~\ref{chap:netstack} 章；\file{/boot} 固件加载已改为内核直接读写 FAT32，见第~\ref{chap:bootfat32} 章。本章硬件、SDIO 和固件下载部分的说明不变，``文件组成''``分层结构''两节已按新文件名更新。

\section{硬件连接}

XR829 WiFi/BT 合封模组挂在 SMHC1 上，原理图参考 \code{100ASK\_T113-Pro\_Base-SCH\_V1.2} p4/p9：

\begin{table}[H]
\centering
\begin{tabular}{lllp{5.6cm}}
\toprule
信号 & 引脚 & 功能号 & 说明 \\
\midrule
SDC1-CLK & PG0 & 2 & 串 33R 电阻 \\
SDC1-CMD & PG1 & 2 & 上拉 \\
SDC1-D0..D3 & PG2..PG5 & 2 & 上拉 \\
CLK-FANOUT1 & PG11 & 5 & 给模组喂 32.768 kHz 时钟 \\
WL\_WAKE\_AP & PG10 & - & 未使用（当前轮询，无中断） \\
WL\_REG\_ON & PG12 & GPIO 输出 & 模组使能，高有效 \\
\bottomrule
\end{tabular}
\end{table}

\begin{itemize}
  \item SMHC1 寄存器基址 \code{0x04021000}（\code{T113\_SMHC\_BASE(1)}），长度 4 KiB；
  \item IRQ 73（GIC SPI 41，\code{T113\_IRQ\_SMHC1}）在平台资源里声明了，但驱动纯轮询，目前没有使用；
  \item 32.768 kHz 由 CCU 的 CLK\_FANOUT1 输出：XR829 没有内部 RTC 时钟，睡眠时钟必须由 SoC 提供，否则模组上电后不响应 CMD5。
\end{itemize}

板级引脚和电源时序描述在 \file{src/include/t113/wifi.h} 的 \code{t113\_mmc\_platform\_data}，具体数值在 \file{board.c} 的 \code{sdc1\_pdata}。

\section{图解：硬件连接与上电时序}

\begin{txtcode}[basicstyle=\ttfamily\scriptsize]
              T113-S3（Cortex-A7）                                    XR829 WiFi/BT 模组
+------------------------------------------+                +----------------------------------+
| SMHC1  @0x04021000（sunxi-smhc.c）        |                | SDIO 从设备                       |
|   CLK  -- PG0  func2 --[33R]-------------+---------------→| SDIO_CLK                          |
|   CMD  -- PG1  func2  上拉 --------------+←--------------→| SDIO_CMD    命令/响应              |
|   D0   -- PG2  func2  上拉 --------------+←--------------→| SDIO_D0     数据（1-bit 只用 D0）   |
|   D1-3 -- PG3..PG5 func2 上拉 -----------+←--------------→| SDIO_D1..D3（4-bit 模式才用）       |
|                                          |                |                                  |
| CCU CLK_FANOUT1 -- PG11 func5 -----------+-- 32.768 kHz -→| LPO 睡眠时钟（模组内无 RTC 晶振）   |
| GPIO PG12 输出 --------------------------+-- WL_REG_ON --→| 电源/复位使能，高有效              |
| GPIO PG10（未用）←-----------------------+-- WL_WAKE_AP --| 唤醒主机（当前全部轮询）            |
| GIC IRQ 73（SPI 41，已声明未用）          |                |                                  |
|                                          |                | 内部 CPU                         |
| SMHC0  @0x04020000 -- PF0..PF5 --→ TF 卡  |                |   AHB 0x08000000  代码/数据 RAM    |
|   分区 1 FAT32：/boot/*_xr829.bin        | -- 固件来源     |   APB 0x09000000  下载 FIFO/控制块 |
|   分区 2 minifs：/                       |                |   PWRCTRL 0x0ac80000              |
+------------------------------------------+                | 2.4 GHz 射频 -- 天线               |
                                                            +----------------------------------+
\end{txtcode}

两个 SD 控制器分工明确：SMHC0 接 TF 卡，负责根文件系统和 \file{/boot}；SMHC1 接 XR829。WiFi 驱动启动时，要先经 SMHC0 从 FAT32 读出三个固件文件，再经 SMHC1 灌进模组。

上电时序在 \code{smhc\_power()}（\code{sunxi-smhc.c}）里，由 \code{sdio\_attach\_card()} 在发 CMD0 之前调用：

\begin{txtcode}[basicstyle=\ttfamily\scriptsize]
                smhc_power(host, 1)                              sdio_attach_card()
                |                                                |
32k FANOUT1  ___+-----------------------------------------------------------------
                 (1) CCU 0xf30 bit4 开 fanout 门控，0xf3c 选 32k 源并开 fanout1，PG11 切 func5
WL_REG_ON    ------+                    +---------------------------------------------
                   |____________________|
                   (2) 拉低             (3) 拉高
                   |<----- 10 ms ------>|<----- power_on_delay_ms = 100 ms ----->|
SDIO_CLK                                                                 400 kHz |||||||
CMD 线                                                                   CMD0 -> CMD5 ...
\end{txtcode}

32k 时钟必须先于 WL\_REG\_ON 拉高：XR829 的内部电源管理靠它计时，没有它模组不会响应 CMD5。\code{PLL\_PERIPH0} 没开时 fanout 也出不来，\code{smhc\_power()} 会打印提示。

\section{文件组成}

\begin{table}[H]
\centering
\begin{tabular}{p{7.0cm}p{7.6cm}}
\toprule
文件 & 角色 \\
\midrule
\code{src/drivers/mmc/host/sunxi-smhc.c} & SMHC（sunxi-mmc）主机控制器驱动：命令 + PIO 数据阶段，platform 总线驱动，轮询 \\
\code{src/drivers/mmc/core.c} & 轻量 MMC/SDIO core：CMD52/CMD53、CCCR/CIS 解析、sdio 总线、卡枚举、\code{mmc\_set\_clock()} \\
\code{src/drivers/net/wireless/xradio/sdio.c} & SDIO function 驱动 \code{xradio\_sdio\_driver}：probe/remove、CMD53 自检、提速 \\
\code{.../xradio/hwio.c} & HIF 寄存器、APB/AHB 间接访问 \\
\code{.../xradio/fwio.c} & DPLL、bootloader 与固件下载、SDD \\
\code{.../xradio/bh.c} & 底半部：HIF 队列读写、序号、固件输入缓冲计数 \\
\code{.../xradio/wsm.c} & WSM 命令编码、confirm、indication 分派、MIB、TX \\
\code{.../xradio/main.c} & 核心 probe：固件 $\rightarrow$ WSM 启动 $\rightarrow$ CONFIGURATION $\rightarrow$ 注册到 mac80211 \\
\code{.../xradio/sta.c} & \code{ieee80211\_ops}：start、JOIN/关联参数、set\_key、poll \\
\code{.../xradio/scan.c}、\code{txrx.c} & 硬件扫描；802.11 帧 $\leftrightarrow$ WSM TX/RX \\
\code{.../xradio/\{xradio,wsm,hwio\}.h} & 驱动私有头文件（WSM wire-format、HIF 定义） \\
\code{src/net/}（core、ethernet、ipv4、wireless、mac80211、wpa、tftp） & 网络栈，见第~\ref{chap:netstack} 章 \\
\code{src/lib/crypto/wpa\_crypto.c}、\code{src/include/crypto/wpa\_crypto.h} & 通用 WPA2-PSK 密码学 \\
\code{src/boards/100ask\_t113/bootpack.c} & 内核侧 bootpack 查询（SPL 把 header 拷到 DDR 固定地址） \\
\code{src/include/t113/bootpack.h} & bootpack 容器格式，与 miniBoot、mkbootpack 共用 \\
\code{src/include/t113/mmc.h} & \code{mmc\_host}/\code{sdio\_card}/\code{sdio\_func}/\code{sdio\_driver}、CMD53 API、CCCR/FBR 定义 \\
\code{src/include/t113/wifi.h} & \code{t113\_mmc\_platform\_data}：sdcN 平台数据（总线引脚 + 模组电源时序） \\
\code{src/boards/100ask\_t113/board.c} & \code{sdc1\_device} 平台设备定义、驱动注册顺序 \\
\code{src/boards/100ask\_t113/driver\_model.c} & 初始化顺序：设备 $\rightarrow$ 驱动 $\rightarrow$ \code{mmc\_rescan\_all()} \\
\code{miniBoot/tools/mkbootpack.c} & 主机侧打包工具：kernel + 固件镜像 $\rightarrow$ bootpack.bin \\
\code{rpi\_bm/Makefile} & 检测 \code{firmware/xr829/*.bin} 并自动填入 \code{BOOTPACK\_EXTRA} \\
\bottomrule
\end{tabular}
\end{table}

固件镜像本身（boot/fw/sdd）是 Tina SDK 的有版权文件，不进入 git（\file{rpi\_bm/.gitignore} 忽略 \file{/firmware/}），构建时从本地目录打入 bootpack。

\section{分层结构}

整体照 Linux 的形状搭（名字和回调都尽量对齐），但实现极小：

\begin{txtcode}[basicstyle=\ttfamily\scriptsize]
platform 总线                          sdio 总线
----------------------                ----------------------
sdc1_device  --match-->  t113_smhc_driver
    |                          | probe: 开 CCU 时钟/复位, pinmux,
    |                          |         复位控制器, 注册 mmc_host
    v                          v
  struct device            struct mmc_host (mmc1)
                                  |  mmc_rescan_all() 枚举
                                  v
                            struct sdio_card
                                  |  sdio_add_funcs()
                                  v
                            struct sdio_func (func 1)
                                  |  device_register() 后作为
                                  |  /bus/sdio/devices/mmc1:0001:1
                                  v
   mmc1:0001:1 --match-->  xradio_sdio_driver ("xradio_wlan")
                                  | probe (sdio.c):
                                  |   使能 function, 设块大小 512
                                  |   CMD53 自检 -> 提速到 12 MHz
                                  | xradio_core_probe (main.c):
                                  |   fwio.c 下载固件 -> WSM startup
                                  |   CONFIGURATION -> ieee80211_register_hw
                                  v
                       struct xradio_common -> mac80211 -> wlan0
\end{txtcode}

代码目录按 Linux 的 bus/core/protocol/MLME 边界组织：

\begin{txtcode}
src/drivers/mmc/
+-- host/sunxi-smhc.c      SoC host controller（硬件层）
+-- core.c                 MMC/SDIO core 与 sdio_bus（总线层）

src/drivers/net/wireless/xradio/
+-- sdio.c                 sdio_driver
+-- hwio.c / fwio.c        HIF 访问 / 固件下载
+-- bh.c / wsm.c           消息传输 / WSM 协议
+-- main.c                 核心 probe，注册到 mac80211
+-- sta.c scan.c txrx.c    ieee80211_ops 实现

src/net/mac80211/ src/net/wireless/ src/net/wpa/   802.11 状态机、cfg80211、supplicant
src/lib/crypto/wpa_crypto.c    与具体网卡无关的 WPA2 primitive
\end{txtcode}

依赖方向保持单向：板级 SMHC $\rightarrow$ SDIO core $\rightarrow$ xradio 驱动 $\rightarrow$ mac80211 $\rightarrow$ cfg80211 $\rightarrow$ supplicant；\code{wpa\_crypto} 只提供算法，不反向依赖驱动。

数据面走的是 SDIO function 1 的 CMD53：

\begin{txtcode}
hwio.c
+-- xradio_reg_read_32/write_32()   CMD53 4 字节, 访问 HIF 寄存器窗口
+-- xradio_apb_write/apb_read_32()  SRAM_BASE_ADDR + SRAM_DPORT 间接访问 APB
+-- xradio_ahb_write_32/read_32()   SRAM_BASE_ADDR + AHB_DPORT 间接访问 AHB
+-- sdio_memcpy_fromio/toio()       CMD53 块/字节传输
        +-- sdio_io_rw_extended()
                +-- mmc_command + mmc_data
                        +-- smhc_request() -> smhc_pio() FIFO 搬运
\end{txtcode}

对应 Linux 里的映射关系：

\begin{table}[H]
\centering
\begin{tabular}{ll}
\toprule
miniOS & Linux / xradio \\
\midrule
\code{mmc\_host} + \code{mmc\_host\_ops} & \code{struct mmc\_host} + \code{mmc\_host\_ops} \\
\code{sdio\_card} / \code{sdio\_func} & \code{mmc\_card} / \code{sdio\_func} \\
\code{sdio\_bus\_type} + \code{sdio\_driver} & \code{sdio\_bus\_type} + \code{sdio\_driver} \\
\code{sdio\_driver\_register()} & \code{sdio\_register\_driver()} \\
\code{sdio\_io\_rw\_extended()} & \code{mmc\_io\_rw\_extended()} \\
\code{mmc\_rescan\_all()} & \code{mmc\_rescan()}（Linux 里在 workqueue 中跑） \\
\code{xradio\_reg\_*} / HIF id & xradio \code{hwio.c} 的 \code{xradio\_reg\_read\_32()} 等 \\
固件下载流程 & xradio \code{fwio.c} 的 \code{xradio\_load\_firmware()} \\
\bottomrule
\end{tabular}
\end{table}

\section{总调用链}

从驱动核心初始化到 \code{wlan0} 注册：

\begin{txtcode}[basicstyle=\ttfamily\scriptsize]
kernel_boot()
+-- t113_driver_core_init()                          boards/100ask_t113/driver_model.c
    +-- platform_bus_init() / i2c_bus_init() / sdio_bus_init()
    +-- t113_board_register_devices()
    |   +-- platform_device_register(&sdc1_device)   只登记，不碰硬件
    +-- t113_board_register_drivers()
    |   +-- platform_driver_register(&t113_smhc_driver)
    |   |   +-- t113_smhc_probe(sdc1)                drivers/mmc/host/sunxi-smhc.c
    |   |       +-- platform_get_resource()          SMHC1 基址 0x04021000
    |   |       +-- CCU 0x84c：复位 → 开门控；0x834：mod clock = HOSC（/2 后 12 MHz）
    |   |       +-- t113_pinmux_apply(sdc1_pins)     PG0..PG5 func 2
    |   |       +-- smhc_reset()；HWRST 拉低 10 us 再拉高
    |   |       +-- NTSR |= NEW_MODE；TMOUT；IMASK = 0（不用中断）
    |   |       +-- smhc_set_ios(400 kHz, 1-bit)
    |   |       +-- mmc_add_host()                   卡扫描推迟
    |   +-- sdio_driver_register(&xradio_sdio_driver)    总线上还没有设备
    +-- mmc_rescan_all()                             不在任何 probe 里，避免持锁注册设备
        +-- sdio_attach_card(host) → sdio_add_funcs(host)   （上一节的状态机）
            +-- device_register(func) → xradio_sdio_probe(func)      xradio/sdio.c
                +-- sdio_enable_func()、sdio_set_block_size(512)
                +-- xradio_cmd53_selftest(verbose = 1)
                |   +-- CMD53 读 HIF_CONFIG 两次，必须一致
                |   +-- CMD53 写/读 HIF_TSET_GEN 三组模式，再恢复原值
                +-- mmc_set_clock(12 MHz) + 自检；失败 → mmc_set_clock(0) 回 400 kHz
                +-- xradio_core_probe(func)                          xradio/main.c
                    +-- xradio_load_firmware(hw_priv)                xradio/fwio.c
                    |   +-- need_image("/boot/boot_xr829.bin")       经 VFS → FAT32 → SMHC0
                    |   +-- need_image("/boot/fw_xr829.bin")
                    |   +-- need_image("/boot/sdd_xr829.bin")        SDD 留在内存给 CONFIGURATION
                    |   +-- CONFIG 的 hw-type 位检查
                    |   +-- parse_sdd() → DPLL；写 TSET_GEN 并回读
                    |   +-- CONTROL.WUP，等 RDY；CONFIG |= ACCESS_MODE
                    |   +-- update_dpllctrl()
                    |   +-- download_bootloader()   AHB 0x08000000，逐字写
                    |   +-- download_firmware()     释放 CPU，握手，1 KiB 块写下载 FIFO
                    |   +-- CONFIG：开 IRQ 位，清 ACCESS_MODE → queue 模式
                    +-- wsm_init(func, 2000)                         xradio/wsm.c
                    |   +-- xradio_bh_init(func)                     xradio/bh.c
                    |   +-- 循环 xradio_bh_poll(10)，直到 STARTUP 0x0801
                    |   |   +-- parse_startup()：num_bufs、buf_size、固件版本
                    |   |       +-- xradio_bh_set_caps()              流控上限
                    |   +-- CONTROL = WUP                            保持唤醒
                    +-- xradio_setup_mac()
                    |   +-- wsm_configuration(MAC, SDD)              CONFIGURATION 0x0009
                    |   +-- wsm_read_mib(CURRENT_TX_POWER)           诊断
                    +-- ieee80211_alloc_hw(&xradio_ops, hw_priv)     net/mac80211/main.c
                    +-- ieee80211_register_hw(hw, mac, xradio_rates, 12)
                        +-- wiphy_new() / wiphy_register()           cfg80211
                        +-- xradio_start()                           ops->start
                        |   +-- OPERATIONAL_POWER_MODE MIB
                        |   +-- wsm_set_edca_default()               4 个 AC 的 EDCA
                        |   +-- wsm_set_tx_policies()                速率重试表
                        +-- register_netdev(wlan0)
                            +-- NETDEV_REGISTER 通知 → wpa_netdev_event()   自动连接
\end{txtcode}

\section{初始化顺序与驱动核心约束}

\code{device\_register()} 会拿驱动核心锁并触发匹配绑定（probe 也在同一把锁下运行）。卡枚举要注册 \code{sdio\_func} 设备，如果放在 SMHC 的 probe 里就会在持锁状态下再次注册设备——死锁。所以 \code{t113\_driver\_core\_init()} 先把所有设备、驱动注册完，最后统一调 \code{mmc\_rescan\_all()}，对应 Linux 把 \code{mmc\_rescan} 扔进 workqueue 的做法。

同理，注册顺序也有要求：\code{xradio\_sdio\_driver} 必须在 \code{mmc\_rescan\_all()} 之前注册，否则 function 设备注册时找不到驱动（本模型没有 deferred-probe 重试；未绑定时会打印 \code{no sdio driver matched}）。

另一个后果是：\textbf{整段固件下载都在 xradio\_sdio\_probe() 里、在驱动核心锁下忙等完成}。这在 Linux 里不可接受（probe 必须快、可睡眠），这里因为只有这一个 SDIO 设备、且下载在 12 MHz 下只需约 0.2 s 量级（失败回退到 400 kHz 时约数秒），所以先保持最简单的同步实现。后续接中断/工作队列时应把固件下载移出 probe。

\code{xradio\_sdio\_probe()} 即使 CMD53 或固件下载失败也返回 0，保持绑定关系，方便驱动模型打印出设备状态并在串口留下诊断信息（若 probe 返回负值，绑定不会建立，现场更难分析）。

\section{图解：SDIO 协议与初始化}

\subsection{命令与响应的格式}

SDIO 沿用 SD 卡的物理层：主机在 CMD 线上发 48 位命令，卡在 CMD 线上回响应，数据走 D0（4-bit 模式下走 D0–D3）。

\begin{txtcode}[basicstyle=\ttfamily\scriptsize]
命令（主机 → 卡，48 位）
 47   46   45 ...... 40   39 .................. 8   7 ... 1   0
+---+----+------------+------------------------+---------+---+
| 0 | 1  | 命令号 6 位 | 参数 32 位              | CRC7    | 1 |
+---+----+------------+------------------------+---------+---+
 起始 方向=主机                                           结束

SMHC 只要软件给出“命令号 + 参数 + 期望的响应类型”，起始位、CRC7 和结束位由硬件生成：
  SMHC_ARG = 参数
  SMHC_CMD = START | 命令号 | RESP_EXPIRE [| LONG_RESP] [| CHK_CRC] [| DATA | WAIT_PRE [| WRITE]]
  轮询 SMHC_RINT 等 CMD_DONE（有数据时再等 DATA_OVER）
  SMHC_RESP0 = 响应里的 32 位内容
\end{txtcode}

SDIO 初始化用到的命令和响应：

\begin{table}[H]
\centering
\small
\begin{tabularx}{\linewidth}{l>{\raggedright\arraybackslash}Xll>{\raggedright\arraybackslash}X}
\toprule
命令 & 名称 & 参数 & 响应 & miniOS 用它做什么 \\
\midrule
CMD0 & GO\_IDLE\_STATE & 0 & 无 & 复位卡的状态机；带 \code{INIT\_SEQ} 位，让 SMHC 先发初始化时钟（规范要求至少 74 个） \\
CMD5 & IO\_SEND\_OP\_COND & 0，之后是 OCR & R4 & 先问电压窗口，再反复协商直到就绪 \\
CMD3 & SEND\_RELATIVE\_ADDR & 0 & R6 & 卡给出 16 位 RCA \\
CMD7 & SELECT\_CARD & RCA << 16 & R1b & 选中卡，进入可传输状态 \\
CMD52 & IO\_RW\_DIRECT & 见下 & R5 & 读写一个字节：CCCR、FBR、CIS \\
CMD53 & IO\_RW\_EXTENDED & 见下 & R5 + 数据 & 批量读写：HIF 寄存器、下载 FIFO、消息队列 \\
\bottomrule
\end{tabularx}
\end{table}

响应里用到的位：

\begin{txtcode}
R4（CMD5，无 CRC）
 31   30..28      27   26..24   23 ............... 0
+---+----------+----+--------+-----------------+
| C | 功能数    | MP | 保留    | OCR 电压窗口     |   C = 1：卡已就绪
+---+----------+----+--------+-----------------+   miniOS 取 OCR & 0x00ff8000（2.7–3.6 V）

R6（CMD3）
 31 ......... 16   15 ......... 0
+--------------+--------------+
| RCA          | 卡状态        |
+--------------+--------------+

R5（CMD52/CMD53）
 31 ... 16   15        14        13..12    11     10   9        8      7 ...... 0
+---------+---------+---------+---------+------+----+--------+------+----------+
| 保留     | CRC 错  | 非法命令 | IO 状态  | 错误 | 0  | 功能号错| 越界  | 读到的数据 |
+---------+---------+---------+---------+------+----+--------+------+----------+
R5_ERROR_MASK = 0xcb00（位 15、14、11、9、8），命中就返回 -EIO
\end{txtcode}

\subsection{初始化状态机}

\begin{txtcode}
sdio_attach_card(host)                            drivers/mmc/core.c
|
+-- smhc_power(host, 1)                           32k + WL_REG_ON 时序
+-- set_ios(400 kHz, 1-bit)
|
+-- CMD0                        卡：任何状态 → Idle
+-- CMD5  arg = 0       → R4    读 OCR 和功能数，不改变状态
|     OCR & 0x00ff8000 == 0 → "card OCR has no 2.7-3.6V window"
+-- CMD5  arg = ocr     → R4    循环，每次间隔 10 ms，最多 100 次
|     R4.C == 1  →  卡：Idle → Ready
+-- CMD3                → R6    RCA；卡：Ready → Stand-by
+-- CMD7  RCA << 16     → R1b   卡：Stand-by → Command（可以收发数据）
|
+-- CMD52 读 CCCR 0x00          CCCR/SDIO 版本
+-- CMD52 读 CCCR 0x08          能力位
+-- CMD52 读 CCCR 0x09..0x0b    公共 CIS 指针
|     +-- parse_cis_manfid()    沿元组链表找 CISTPL_MANFID(0x20)
|                               → vendor 0x0a9e，device 0x2282
+-- sdio_add_funcs(host)        对每个 function n = 1..功能数
      +-- CMD52 读 FBR n×0x100 + 0x00     标准接口类别
      +-- CMD52 读 FBR n×0x100 + 0x09     本 function 的 CIS 指针 → MANFID
      +-- device_register(&func->dev)     名字 "mmc1:0001:1"
            +-- sdio_bus_match()          按 xradio_sdio_ids[] 匹配
                  +-- xradio_sdio_probe(func)            xradio/sdio.c
                        +-- sdio_enable_func()
                        |   +-- CMD52 读改写 CCCR 0x02 IOEx |= 1 << 1
                        |   +-- CMD52 轮询 CCCR 0x03 IORx bit1   function 1 就绪
                        +-- sdio_set_block_size(512)
                        |   +-- CMD52 写 FBR1 + 0x10 = 0x00，+ 0x11 = 0x02
                        +-- xradio_cmd53_selftest(verbose)      CMD53 能不能用
                        +-- mmc_set_clock(12 MHz) → 再自检一次，不行退回 400 kHz
                        +-- xradio_core_probe(func)            进入固件阶段
\end{txtcode}

miniOS 没有发 CMD52 去写 CCCR 0x07（总线宽度），所以卡一直停在默认的 1-bit模式；也没有用 CMD52 写 CCCR 0x04（IENx）打开卡中断。这两项是 Linux 的\code{mmc\_attach\_sdio()} 会做、而 miniOS 刻意省掉的。

\section{SDIO 枚举}

\subsection{CMD52：单字节读写}

所有 CCCR/FBR 寄存器访问走 \code{sdio\_io\_rw\_direct()}，命令参数位域：

\begin{txtcode}
bit 31     R/W        1 = 写
bit 30:28  function   0 = CCCR/FBR 空间, 1..7 = 各 function 寄存器
bit 27     RAW        写后立即回读
bit 26     保留 0
bit 25:9   地址       17 位
bit  7:0   数据       写方向是数据, 读方向是 stuff
\end{txtcode}

响应 R5 的错误位掩码取 \code{0xcb00}（与 Linux \code{mmc\_io\_rw\_direct()} 相同）：

\begin{txtcode}
bit 15     COM_CRC_ERROR     CRC 校验失败
bit 14     ILLEGAL_COMMAND
bit 13:12  IO_CURRENT_STATE  卡的 I/O 状态，不算错误
bit 11     ERROR             通用错误
bit  9     FUNCTION_NUMBER   function 号超范围
bit  8     OUT_OF_RANGE      地址超范围
\end{txtcode}

封装（\file{drivers/mmc/core.c}）：

\begin{itemize}
  \item \code{sdio\_readb/writeb}：指定 function 的单字节读写；
  \item \code{sdio\_f0\_readb/writeb}：固定 function 0，即 CCCR/FBR 配置空间；
  \item \code{sdio\_enable\_func}：IOEx 置位 function 位后轮询 IORx 对应位，最长 1 秒；
  \item \code{sdio\_set\_block\_size}：写 FBR+0x10/0x11，并同步记录到 \code{sdio\_func.blksz}（CMD53 块模式要用）。
\end{itemize}

\subsection{CCCR / FBR / CIS}

CCCR（function 0 公共配置）：

\begin{table}[H]
\centering
\begin{tabular}{llp{7.2cm}}
\toprule
地址 & 名称 & 作用 \\
\midrule
0x00 & CCCR & 低 4 位 CCCR 版本, 高 4 位 SD 版本 \\
0x02 & IOEx & 各 function I/O 使能（写） \\
0x03 & IORx & 各 function 就绪状态（只读） \\
0x08 & CAPS & 能力（低功耗、DMA、高速支持） \\
0x09..0x0b & CIS & 指向公共 CIS 元组链表的 24 位指针（小端） \\
\bottomrule
\end{tabular}
\end{table}

FBR（每个 function n 一份，基址 \code{n*0x100}）：

\begin{table}[H]
\centering
\begin{tabular}{llp{7.6cm}}
\toprule
偏移 & 名称 & 作用 \\
\midrule
+0x00 & STD\_IF & 低 4 位 standard interface code；XR829 不用它匹配（驱动按 CIS 的 vendor/device 匹配） \\
+0x09.. & CIS & 该 function 自己的 CIS 指针 \\
+0x10/0x11 & BLKSIZE & 块大小低/高字节 \\
\bottomrule
\end{tabular}
\end{table}

CIS 是元组链表：每项一字节 code + 一字节 link（长度），跟着内容。\code{parse\_cis\_manfid()} 沿链表走，找 \code{CISTPL\_MANFID(0x20)}，取前 4 字节得到 vendor/device。

实测这块 100ASK T113-Pro 上的 XR829 报 \textbf{0x0a9e:0x2282}；开源 xradio 驱动里的老 XR819 是 0x0020:0x2281（ST-Ericsson 历史遗留 ID）。\code{xradio\_sdio\_ids[]} 两个都匹配，另外各留一条只限 vendor 的兜底项：

\begin{ccode}
{ SDIO_ANY_ID, 0x0a9e, 0x2282 },        /* XR829, 本板实测 */
{ SDIO_ANY_ID, 0x0020, 0x2281 },        /* XR819 per xradio */
{ SDIO_ANY_ID, 0x0a9e, SDIO_ANY_ID },
{ SDIO_ANY_ID, 0x0020, SDIO_ANY_ID },
{ 0, 0, 0 },
\end{ccode}

\section{CMD53 与数据通路}

\code{sdio\_io\_rw\_extended()} 实现 CMD53，参数位域：

\begin{txtcode}
bit 31     R/W        1 = 写
bit 30:28  function   function 1
bit 27     block     1 = 块模式, 0 = 字节模式
bit 26     incr      1 = 地址自增（内存型）, 0 = 固定地址（FIFO 口）
bit 25:9   地址       17 位
bit  8:0   count      字节模式 = 字节数（512 编码为 0）
                      块模式 = 块数
\end{txtcode}

约束与选择逻辑（都强制 4 字节对齐）：

\begin{table}[H]
\centering
\begin{tabular}{llp{7.4cm}}
\toprule
条件 & 模式 & 说明 \\
\midrule
\code{len <= 512} & 字节模式 & 单次传输，\code{data.blksz = len}, \code{blocks = 1} \\
\code{len > 512} & 块模式 & \code{len} 必须是 \code{f->blksz} 的整数倍，块数 $\leq$ 511 \\
\bottomrule
\end{tabular}
\end{table}

对外接口：

\begin{itemize}
  \item \code{sdio\_memcpy\_fromio(func, dst, addr, len)} / \code{sdio\_memcpy\_toio()}：地址自增的 CMD53，HIF 的寄存器、FIFO、队列都走它；
  \item HIF 寄存器访问封装成 \code{xradio\_reg\_read\_32/write\_32()}（4 字节，hwio.c），声明在 \code{mmc.h} 末尾；
  \item 1 KiB 固件块 = 2 $\times$ 512 字节块传输。
\end{itemize}

\subsection{SMHC PIO 数据阶段}

数据阶段不用 DMA、不用中断，CPU 直接搬 FIFO：

\begin{txtcode}
smhc_request() 带 mmc_data 时：
+-- CMD 寄存器置 DATA(b9) | WAIT_PRE(b13) | WRITE(b10)
+-- GCTRL |= AHB(b31) | FIFO_RST(b1)，等 FIFO_RST 自清
+-- BLKSZ = blksz, BYTECNT = blksz * blocks
+-- 启动命令后进 smhc_pio()
|   +-- 写：STATUS.FIFO_FULL(b3) 为 0 时写 SMHC_FIFO(0x200)
|   +-- 读：STATUS.FIFO_EMPTY(b2) 为 0 时读 SMHC_FIFO
|   +-- 每搬 4 字节检查一次 RINT 错误位，有错立即 -EIO
|   +-- 连续 ~200 ms 没有进展 -> -ETIMEDOUT
+-- 等 RINT.CMD_DONE(b2) | DATA_OVER(b3)
\end{txtcode}

\code{SMHC\_FIFO} 在偏移 \code{0x200}（D1/T113 的 NCAT2 布局）。出错时和纯命令路径一样：\code{smhc\_reset()} + 重新应用时钟，避免状态机卡死。

\section{SMHC 主机控制器}

\file{sunxi-smhc.c} 是命令级轮询驱动：命令 + PIO 数据，不接中断。

用到的寄存器（偏移）：

\begin{table}[H]
\centering
\begin{tabular}{llp{8.6cm}}
\toprule
偏移 & 名称 & 用途 \\
\midrule
0x000 & GCTRL & 软/FIFO/DMA 复位（\code{0x7}），bit1 FIFO 复位，bit31 CPU 访存 FIFO \\
0x004 & CLKCR & bit16 时钟使能, bit17 省电, bit31 MASK\_D0, [7:0] 分频 \\
0x008 & TMOUT & 超时配置（初始化写 \code{0xffffff40}） \\
0x00c & WIDTH & 0 = 1-bit, 1 = 4-bit \\
0x010 & BLKSZ & 数据阶段块大小 \\
0x014 & BYTECNT & 数据阶段总字节数 \\
0x018 & CMD & bit31 启动, b21 UPCLK\_ONLY, b15 INIT\_SEQ, b13 WAIT\_PRE, b10 WRITE, b9 DATA, b8/7/6 CRC/长响应/响应 \\
0x01c & ARG & 命令参数 \\
0x020 & RESP0 & 32 位响应（136 位响应在 RESP0..3） \\
0x030 & IMASK & 全 0：不用中断 \\
0x038 & RINT & b2 命令完成, b3 数据完成, b8 响应超时, 错误掩码 0xbfc2 \\
0x03c & STATUS & b2 FIFO 空, b3 FIFO 满, b9 卡忙（R1b 判忙用） \\
0x05c & NTSR & bit31 new timing 模式 \\
0x078 & HWRST & 0 = 复位, 1 = 释放 \\
0x200 & FIFO & PIO 数据端口（NCAT2） \\
\bottomrule
\end{tabular}
\end{table}

时钟（CCU，\code{ccu-sun20i-d1} 布局）：

\begin{itemize}
  \item 模块时钟 \code{0x830 + 4N}：bit31 gate，bit24-26 mux（0 = HOSC 24 MHz），P b8-9、M b0-3，\textbf{固定后分频 2}，即有效源时钟 12 MHz；
  \item 总线 gate/复位 \code{0x84c}：gate bit N，reset bit 16+N；
  \item fanout gate \code{0xf30} bit4、fanout 配置 \code{0xf3c}：每通道 3 位源选择（0 = 32k）+ bit21+n gate。
\end{itemize}

\code{smhc\_set\_ios()} 的换频序列是 D1/T113 的坑点，Linux sunxi-mmc 同款：

\begin{enumerate}
  \item 关 CLKCR\_ENABLE、置 MASK\_D0，发一次 UPCLK\_ONLY 命令；
  \item 写新分频 + ENABLE + MASK\_D0，再发 UPCLK\_ONLY；
  \item 清掉 MASK\_D0。
\end{enumerate}

换频期间必须把 DATA0 屏蔽掉，否则卡时钟抖动会让总线上出现非法电平。分频公式 \code{div = ceil(12M / (2*hz))} 保证实际时钟只会更低：

\begin{table}[H]
\centering
\begin{tabular}{lll}
\toprule
目标 & 分频 & 实际时钟 \\
\midrule
400 kHz（识别/下载回退） & 15 & 400 kHz \\
12 MHz（源时钟上限，提速目标） & 0 & 12 MHz \\
\bottomrule
\end{tabular}
\end{table}

XR829 驱动用 \code{mmc\_set\_clock(host, 12 MHz)} 提速，\code{mmc\_set\_clock(host, 0)} 回退到 host 的 \code{f\_init}（400 kHz）。12 MHz 是本板 SMHC 源时钟（HOSC 24 MHz 固定 /2）能给出的上限，配合 1-bit 总线；没有做 high-speed 模式（那需要 CMD6 切时序）。

命令超时/错误处理：轮询 RINT，命中错误掩码就按 RESP\_TMO 区分 \code{-ETIMEDOUT} / \code{-EIO}；出错后控制器状态机可能卡死，所以每次错误都做 \code{smhc\_reset()} + 重设时钟再返回。

\section{HIF 寄存器与间接寻址}

XR829 的 host 接口（HIF）固定在 SDIO function 1 的地址空间，遵循 xradio \code{hwio.h} 的寻址方式：

\begin{txtcode}
SDIO 地址 = (buf_id << 7) | (reg_id << 2)
\end{txtcode}

即 32 个 128 字节的 buffer 窗口；寄存器 id 在窗口内按 4 字节步进（本驱动用到 0..6）。direct 模式用 buf\_id 0 访问寄存器：

\begin{table}[H]
\centering
\begin{tabular}{llp{8.2cm}}
\toprule
reg\_id & 名称 & 作用 \\
\midrule
0 & HIF\_CONFIG & 全局配置：访问模式、prefetch、CPU 复位/时钟、IRQ 使能、hw-type \\
1 & HIF\_CONTROL & WUP/RDY 握手、NEXT\_LEN（待收消息长度） \\
2 & HIF\_IN\_OUT\_QUEUE & queue 模式的数据窗口 \\
3 & HIF\_AHB\_DPORT & AHB 间接访问端口 \\
4 & HIF\_SRAM\_BASE\_ADDR & 间接访问的目标地址 \\
5 & HIF\_SRAM\_DPORT & APB/SRAM 间接访问端口 \\
6 & HIF\_TSET\_GEN\_R\_W & DPLL 配置寄存器（写完由硬件读取） \\
\bottomrule
\end{tabular}
\end{table}

CONFIG / CONTROL 关键位：

\begin{txtcode}
CONFIG b10 ACCESS_MODE      1 = direct（端口）模式, 0 = queue（消息）模式
CONFIG b11 AHB_PFETCH       AHB 读预取中
CONFIG b12 CPU_CLK_DIS      CPU 时钟关
CONFIG b13 PFETCH           SRAM 读预取中
CONFIG b14 CPU_RESET        CPU 复位
CONFIG b16|b17 IRQ 使能
CONFIG b26 hw-type 标识     固件下载前必须为 1
CONTROL b11:0 NEXT_LEN      下一条消息长度（16-bit 字为单位）
CONTROL b12 WUP             写 0 唤醒设备
CONTROL b13 RDY             设备就绪
\end{txtcode}

设备内存不在 SDIO 地址空间里，要通过间接窗口访问：

\begin{txtcode}
写:  SRAM_BASE_ADDR = 目标地址
     CMD53 写 SRAM_DPORT（APB 空间）或 AHB_DPORT（AHB 空间）

读:  SRAM_BASE_ADDR = 目标地址
     CONFIG |= PFETCH（SRAM）或 AHB_PFETCH（AHB）
     轮询 CONFIG 直到该位清零（退避等待, 最多约 190 ms）
     CMD53 读对应端口
\end{txtcode}

地址空间：

\begin{itemize}
  \item AHB \code{0x08000000} 起：设备 CPU 的内存，bootloader 镜像写到这里；
  \item APB \code{0x09000000} 起：设备外设空间，下载 FIFO/控制块在这里；
  \item AHB \code{0x0ac80000}：WLAN 电源控制（PWRCTRL）寄存器组。
\end{itemize}

所有 HIF 传输都是地址自增的 CMD53（与 xradio 的 hwio 一致），寄存器访问是 4 字节字节模式传输。

\section{固件下载}

\subsection{镜像}

\begin{table}[H]
\centering
\begin{tabular}{llllp{4.6cm}}
\toprule
bootpack 名称 & 文件 & 大小 & 加载地址 & 作用 \\
\midrule
\code{boot\_xr829} & boot\_xr829.bin & 2184 B & 0x40100000 & 设备侧下载器（写 AHB 内存后由 CPU 执行） \\
\code{fw\_xr829} & fw\_xr829.bin & 174624 B & 0x40110000 & 主固件，经下载 FIFO 灌入 \\
\code{sdd\_xr829} & sdd\_xr829.bin & 744 B & 0x40140000 & RF/时钟校准数据，取参考时钟频率 \\
\bottomrule
\end{tabular}
\end{table}

目录里另有 \code{etf\_xr829.bin}、\code{fw\_xr829\_bt.bin}（蓝牙固件），当前不打包。

\subsection{流程}

\code{xradio\_load\_firmware()} 按 xradio \code{fwio.c xradio\_load\_firmware()} 的顺序执行：

\begin{enumerate}
  \item \code{need\_image()} 经 VFS 从 \file{/boot} 读三个镜像（正常是 FAT32；FAT32 挂载失败时是 bootfs 的 DDR 视图）。缺任何一个就打印 \code{cannot load <路径>} 并放弃，\code{xradio\_sdio\_probe()} 仍返回 0，保持绑定。
  \item CMD53 读 CONFIG，检查 hw-type 位（bit26）：不是 XRadio HIF 就退出。
  \item 解析 SDD：遍历 \code{\{u8 id, u8 len, data\}} 元素，取 \code{0xc5}（reference frequency，kHz，小端 16 位），按表换算 DPLL 值：

  \begin{table}[H]
  \centering
  \begin{tabular}{ll}
  \toprule
  参考时钟 & DPLL 值 \\
  \midrule
  13 MHz & 0x1d89d241 \\
  16 MHz & 0x000001e1 \\
  16.8 MHz & 0x124931c1 \\
  19.2 MHz & 0x00000191 \\
  24 MHz & 0x00000141 \\
  26 MHz & 0x0ec4f121 \\
  33.6 MHz & 0x092490e1 \\
  38.4 MHz & 0x100010c1 \\
  40 MHz & 0x000000c1 \\
  48 MHz & 0x000000a1 \\
  52 MHz & 0x07627091 \\
  \bottomrule
  \end{tabular}
  \end{table}

  \item CMD53 把 DPLL 写进 TSET\_GEN\_R\_W，延迟 5 ms 后回读校验。
  \item CONTROL.WUP 写 0 唤醒设备，退避轮询 CONTROL.RDY（累计约 0.6 s）。
  \item CONFIG 置 ACCESS\_MODE（direct 端口模式）。
  \item \code{update\_dpllctrl()}：在 AHB 0x0ac80000 的 PWRCTRL 上等 WLAN 仲裁（ARBT 且非 BT），把 DPLL 值写进 DPLL\_CTRL（带使能 bit0），再置 COMMON\_CFG.DONE；最后清 START\_CFG。
  \item \code{download\_bootloader()}：把 boot\_xr829 以 32 位字逐个写入 AHB 0x08000000（每字一次 CMD53 AHB 间接写）。
  \item 清 CONFIG.CPU\_RESET 和 CPU\_CLK\_DIS，释放设备 CPU。
  \item 轮询 APB 的 IMAGE\_SIZE 寄存器直到读到 \code{I\_AM\_HERE}（100 $\times$ 10 ms）：设备侧 downloader 已跑起来。
\end{enumerate}

然后 \code{download\_firmware()}：

\begin{enumerate}
  \item 初始化控制块：IMAGE\_SIZE = \code{ARE\_YOU\_HERE}，PUT/GET = 0，STATUS = \code{PENDING}，FLAGS = 0。
  \item 再次检查块大小：1024 必须能被 \code{f->blksz} 整除且 \code{blksz $\leq$ 1024}，否则设成 512。
  \item IMAGE\_SIZE = 实际镜像字节数。
  \item 以 1 KiB 为块循环（镜像 171 块）：
    \begin{itemize}
      \item 读 STATUS，必须仍是 PENDING；
      \item 数据拷进 1 KiB 对齐缓冲，最后一块补零；
      \item 等 FIFO 空间：\code{put - get <= 32 KiB - 1 KiB}（留一块余量）；
      \item CMD53 把 1 KiB 写到 FIFO 环内偏移 \code{put \& 0x7fff}；
      \item \code{put += 实际字节数}，写回 PUT 寄存器；
      \item 每 32 块打印一次进度（32/64/96/128/160 KiB）。
    \end{itemize}
  \item 退避轮询 STATUS 直到非 PENDING；必须等于 \code{SUCCESS}（0），否则打印 STATUS 和 TRACE\_PC。
  \item CONFIG 使能设备中断位（b16|b17），再清 ACCESS\_MODE 切到 queue（消息）模式，延迟 10 ms。
\end{enumerate}

\subsection{下载 FIFO 协议}

\begin{txtcode}
APB 0x09004000  32 KiB 下载 FIFO（put/get 环形窗口, 掩码 0x7fff）
APB 0x0900ff80  控制块:
  +0x00 IMAGE_SIZE    握手: ARE_YOU_HERE -> I_AM_HERE, 随后写镜像大小
  +0x04 FLAGS
  +0x08 PUT           生产者字节计数（host 更新）
  +0x0c TRACE_PC      失败时设备侧 PC
  +0x10 GET           消费者字节计数（设备更新）
  +0x14 STATUS        PENDING / SUCCESS / 错误码
\end{txtcode}

常量：\code{ARE\_YOU\_HERE=0x87654321}、\code{I\_AM\_HERE=0x12345678}、\code{PENDING=0xffffffff}、\code{SUCCESS=0}。

这是典型的单生产者/单消费者环形队列：host 只推进 PUT，设备只推进 GET；host 靠 \code{put - get} 判断剩余空间。最后一块补零写入但 PUT 只按实际镜像长度推进，设备按 IMAGE\_SIZE 截断。

\subsection{进入 queue 模式}

固件跑起来后，数据通路从``端口 + 间接地址''切换为``消息队列''：

\begin{itemize}
  \item 读：\code{HIF\_IN\_OUT\_QUEUE} 端口，RX buffer id 从 1 开始按 \code{+1 \& 3} 轮转（\code{buf\_id\_rx}）；
  \item 写：同一端口用 TX buffer id（\code{buf\_id\_tx} 在 0..31 间轮转，见“通信协议”一节）；
  \item 消息长度由 CONTROL.NEXT\_LEN 给出。
\end{itemize}

\section{通信协议：从 SDIO 到 WSM}

\subsection{协议分层}

\begin{txtcode}[basicstyle=\ttfamily\scriptsize]
+-----------------------------------------------------------------+
| 802.11 帧（mac80211 组好的管理帧、数据帧）                          |  net/mac80211
+-----------------------------------------------------------------+
| WSM 消息：{u16 len, u16 id} + 负载                                 |  xradio/wsm.c
|   请求 0x00xx → confirm 0x04xx；固件主动发 indication 0x08xx        |
+-----------------------------------------------------------------+
| HIF 队列：IN_OUT_QUEUE 端口、buf_id 轮转、3 位序号、NEXT_LEN、       |  xradio/bh.c
|           固件输入缓冲计数                                          |
+-----------------------------------------------------------------+
| HIF 寄存器窗口：SDIO 地址 = HIF_ADDR(buf_id, reg_id)                |  xradio/hwio.c
+-----------------------------------------------------------------+
| SDIO：CMD52 一个字节（function 0 配置空间）                          |  drivers/mmc/core.c
|       CMD53 多字节（function 1，字节模式 ≤ 512 B，或 512 B 块模式）  |
+-----------------------------------------------------------------+
| SMHC1：命令寄存器 + FIFO，CPU 轮询搬运                               |  mmc/host/sunxi-smhc.c
+-----------------------------------------------------------------+
| CLK / CMD / D0                                                    |
+-----------------------------------------------------------------+
\end{txtcode}

XR829 的 HIF 有两种工作模式，\code{CONFIG.ACCESS\_MODE}（bit10）切换：

\begin{table}[H]
\centering
\small
\begin{tabularx}{\linewidth}{>{\raggedright\arraybackslash}Xl>{\raggedright\arraybackslash}X}
\toprule
模式 & 何时 & 怎么访问设备 \\
\midrule
direct（ACCESS\_MODE = 1） & 固件下载期间 & 寄存器窗口 + \code{SRAM\_BASE\_ADDR} 间接访问 AHB/APB 内存 \\
queue（ACCESS\_MODE = 0） & 固件跑起来以后 & 只通过 \code{IN\_OUT\_QUEUE} 端口收发 WSM 消息 \\
\bottomrule
\end{tabularx}
\end{table}

\subsection{HIF 地址}

\begin{txtcode}
CMD53 参数里的 17 位地址（function 1）
 16 ... 12   11 ...... 7    6 ...... 2    1  0
+----------+-------------+-----------+------+
|    0     | buf_id 5 位  | reg_id    | 0  0 |
+----------+-------------+-----------+------+
HIF_ADDR(buf_id, reg_id) = (buf_id & 0x1f) << 7 | reg_id << 2

寄存器访问：buf_id = 0，reg_id = 0..6（CONFIG、CONTROL、IN_OUT_QUEUE、AHB_DPORT、
                                     SRAM_BASE_ADDR、SRAM_DPORT、TSET_GEN）
队列写：    buf_id = tx 序号 0..31，reg_id = 2      每写一条消息 +1
队列读：    buf_id = rx 序号 + 1（1..4），reg_id = 2  每读一次 +1，模 4
\end{txtcode}

\subsection{WSM 消息格式}

\begin{txtcode}
字节  0 ...... 1      2 ............... 3      4 ...
    +-------------+----------------------+--------------+
    | len（LE）    | id（LE）              | 负载           |   len 含这 4 字节头
    +-------------+----------------------+--------------+

id 的位
 15 ... 13   12   11   10    9 ...... 6     5 ...... 0
+---------+----+----+----+-----------+--------------+
| 序号     |保留| IND| CFM| 接口/链路号| 消息号        |
+---------+----+----+----+-----------+--------------+
WSM_MSG_ID(id) = id & 0x0c3f     去掉序号和链路号
0x00xx 请求    0x04xx confirm（= 请求 | 0x0400）   0x08xx indication
\end{txtcode}

TX 请求（\code{wsm\_tx\_frame()}，0x0004）在 4 字节头之后还有 20 字节，一共 24 字节：

\begin{table}[H]
\centering
\small
\begin{tabularx}{\linewidth}{l>{\raggedright\arraybackslash}X>{\raggedright\arraybackslash}X}
\toprule
偏移 & 字段 & miniOS 填什么 \\
\midrule
4 & packetID（u32） & mac80211 的 cookie，TX confirm 原样带回 \\
8 & queueId & 链路号左移 2 位再加队列号；STA 的链路号是 0，队列 VO 3、VI 2、BE 0、BK 1 \\
9 & more & 0 \\
10 & DataOffset & 24，802.11 帧从这里开始 \\
11 & EPTA 优先级 & 0 \\
12 & TxRateEntry（u16） & 管理帧 0x000f，数据帧 0x00a4 \\
14 & TxPolicyIndex & 管理帧 0，数据帧 1 \\
15 & ExpireTimeSetting & 0 \\
16 & expireTime（u32）+ 保留 & 0，用固件默认 \\
24 & 802.11 帧 & 头 + 数据，不带 FCS \\
\bottomrule
\end{tabularx}
\end{table}

RX indication（0x0804）的 16 字节头：

\begin{table}[H]
\centering
\small
\begin{tabularx}{\linewidth}{l>{\raggedright\arraybackslash}X}
\toprule
偏移（相对负载） & 字段 \\
\midrule
0 & status（u32），非 0 表示出错或没有钥匙，驱动丢弃 \\
4 & channel（u16） \\
6 & 收到的速率表项（u16） \\
11 & RCPI/RSSI \\
12 & flags（u32），低 3 位是加密类型：0 无、2 TKIP、3 AES \\
16 & 802.11 帧（固件已解密，但 IV 和 MIC 还在） \\
\bottomrule
\end{tabularx}
\end{table}

用到的消息：

\begin{table}[H]
\centering
\small
\begin{tabularx}{\linewidth}{l>{\raggedright\arraybackslash}X>{\raggedright\arraybackslash}X>{\raggedright\arraybackslash}X}
\toprule
方向 & id & 名字 & 谁处理 \\
\midrule
请求 & 0x0004 & TX & \code{wsm\_tx\_frame()} \\
请求 & 0x0005 / 0x0006 & READ\_MIB / WRITE\_MIB & \code{wsm\_read\_mib()} / \code{wsm\_write\_mib()} \\
请求 & 0x0007 & START\_SCAN & \code{wsm\_scan()} \\
请求 & 0x0009 & CONFIGURATION（带 SDD） & \code{wsm\_configuration()} \\
请求 & 0x000b & JOIN & \code{wsm\_join()} \\
请求 & 0x000c / 0x000d & ADD\_KEY / REMOVE\_KEY & \code{wsm\_add\_key\_*()} / \code{wsm\_remove\_key()} \\
请求 & 0x0010 / 0x0011 / 0x0013 & SET\_PM / SET\_BSS\_PARAMS / SET\_EDCA & \code{sta.c} \\
confirm & 0x0404 / 0x041e & TX confirm / 多个 TX confirm & \code{xradio\_tx\_confirm\_cb()} \\
indication & 0x0800 & EXCEPTION & 打印并 \code{xradio\_bh\_halt()} \\
indication & 0x0801 & STARTUP & \code{parse\_startup()} \\
indication & 0x0804 & RX & \code{xradio\_rx\_cb()} \\
indication & 0x0805 & EVENT（BSS\_LOST 等） & \code{xradio\_event\_cb()} \\
indication & 0x0806 & SCAN\_COMPLETE & \code{xradio\_scan\_complete\_cb()} \\
indication & 0x0809 & SET\_PM 完成 & 记标志 \\
\bottomrule
\end{tabularx}
\end{table}

\subsection{一次队列传输}

\begin{txtcode}
读（rx_one）
+-- 有上次捎带的 CONTROL 就用它，否则 CMD53 读 HIF_CONTROL（4 字节）
+-- len = NEXT_LEN × 2（NEXT_LEN 以 16 位字为单位）；为 0 表示没有消息
+-- rd = align_len(len + 2)          ≤ 512 → 4 字节对齐的字节模式
|                                     > 512 → 512 B 的整数块
+-- CMD53 读 HIF_ADDR(buf_id_rx + 1, IN_OUT_QUEUE)，rd 字节
+-- 末尾 2 字节 = 下一个 CONTROL 值（捎带，省一次寄存器读）
+-- 校验 3 位序号（期望 rx_seq，之后 +1 模 8）
+-- 一次传输里可能有多条消息（multi-rx），逐条：
    +-- release_buffers()            confirm 归还固件输入缓冲
    +-- wsm_handle_rx()

写（xradio_bh_tx）
+-- while hw_bufs_used ≥ num_bufs：xradio_bh_poll(10)   等 confirm 还缓冲，最多 2 s
+-- 写 len，把 tx_seq 填进 id 的 15..13 位
+-- wlen = align_len(len)，短消息至少 16 字节，补零
+-- CMD53 写 HIF_ADDR(buf_id_tx, IN_OUT_QUEUE)
+-- buf_id_tx = (buf_id_tx + 1) & 31；tx_seq++；hw_bufs_used++
\end{txtcode}

\subsection{流控：固件输入缓冲}

STARTUP indication 告诉主机固件有几个输入缓冲（\code{num\_bufs}）、每个多大（\code{buf\_size}）。每个请求（包括 TX）占一个，每个 confirm 还一个：

\begin{txtcode}
主机 hw_bufs_used                          固件输入缓冲（num_bufs 个）
0 -- wsm_cmd(JOIN) - CMD53 写 -----------→ [J][ ][ ][ ]...        1
1 -- wsm_tx_frame(EAPOL M2) -------------→ [J][T][ ][ ]...        2
2 ←- JOIN confirm 0x040b - CMD53 读 ------ [ ][T][ ][ ]...        1   release_buffers()
1 ←- TX confirm 0x0404（packetID）-------- [ ][ ][ ][ ]...        0
满 → xradio_bh_tx() 先轮询收 confirm，腾出缓冲再写
\end{txtcode}

\subsection{命令和 confirm 的配对}

\begin{txtcode}
wsm_cmd(req, payload, len, cfm, cap, timeout)
+-- mutex_lock(&command_mutex)                一次只有一个命令在等
+-- wait_cfm = req | 0x0400，cfm_done = 0
+-- xradio_bh_tx(cmd_buf)
+-- while !cfm_done：xradio_bh_poll(5)         自己驱动底半部
|   +-- rx_one → wsm_handle_rx
|       +-- id == wait_cfm → 拷负载，cfm_status = 负载前 4 字节，cfm_done = 1
|       +-- TX confirm → ieee80211_tx_status()   照常分派，不影响等待
|       +-- indication → 各自回调（只入队或记标志）
+-- 返回 cfm_status（0 = WSM_STATUS_SUCCESS）
\end{txtcode}

Linux 的 xradio 用 \code{wait\_event\_timeout()} 睡眠，由 bh 线程收到 confirm 后唤醒；miniOS 没有中断也没有 bh 线程，等 confirm 的人自己调\code{xradio\_bh\_poll()}。为防止回调里再次进入底半部，\code{bh.in\_rx} 会让嵌套的\code{xradio\_bh\_poll()} 返回 \code{-EBUSY}。

\subsection{启动握手}

固件下载完、切到 queue 模式后，固件先主动发 STARTUP indication。驱动在它之前不能发任何命令：

\begin{txtcode}
xradio_core_probe()
+-- xradio_load_firmware()        结束时 ACCESS_MODE = 0（queue 模式）
+-- wsm_init(func, 2000)
|   +-- xradio_bh_init()          num_bufs = 1，序号和 buf_id 清零
|   +-- 循环 xradio_bh_poll(10)
|       +-- wsm_handle_rx(0x0801) → parse_startup()
|           +-- num_bufs、buf_size → xradio_bh_set_caps()
|           +-- hw_id、fw_type、fw_api、固件标签
|           +-- status == 0 → wsm.ready = 1
+-- CONTROL = WUP                 一直保持唤醒（不做省电）
+-- wsm_configuration(MAC, SDD)   第一条命令必须是 CONFIGURATION：把 SDD 交给固件
    +-- confirm 带回固件实际使用的 MAC、频段和速率掩码
\end{txtcode}

\section{固件打包与加载（bootpack，现为后备路径）}

\textbf{注}：现在固件的正常来源是内核直接读写 SD 卡 FAT32 分区上的 \file{/boot/*.bin}（见第~\ref{chap:bootfat32} 章）；本节的 bootpack 预载 + DDR 交接是 FAT32 挂载失败时的后备路径。

固件比 SRAM 大得多，必须放在 DDR 里由 SPL 预加载。miniOS 用一个小容器把 kernel 和附加镜像一起放到 SD 卡：

\begin{txtcode}
SD 8 KiB      Boot0（miniBoot SPL）
SD 128 KiB    bootpack: 512 B header + 镜像, 每项 512 对齐
\end{txtcode}

bootpack header（\file{src/include/t113/bootpack.h}，miniBoot、miniOS、mkbootpack 三方共用）：

\begin{table}[H]
\centering
\begin{tabular}{ll}
\toprule
字段 & 说明 \\
\midrule
magic / version & \code{"MPAK"} / 1 \\
count / total\_size & 条目数（$\leq$8）/ 总字节数 \\
entry[] & \code{name[16]}、offset、size、load\_addr、crc32 \\
header\_crc32 & 除本字段外整个 header 的 CRC32 \\
\bottomrule
\end{tabular}
\end{table}

SPL 按条目把镜像拷到 \code{load\_addr} 并校验 CRC32，再把 header 复制到 \code{BOOTPACK\_DDR\_ADDR = 0x401ff000}，最后跳到 kernel。内核通过 \code{t113\_bootpack\_find(name)} 查条目，启动时 \code{t113\_bootpack\_dump()} 打印全部条目。

DDR 布局（AUX 窗口是 kernel 下方的 2 MiB）：

\begin{txtcode}
0x40000000  AUX_BASE
0x40100000  boot_xr829.bin   (2184 B)
0x40110000  fw_xr829.bin     (174624 B, 结束于 0x4013aa20)
0x40140000  sdd_xr829.bin    (744 B)
0x401ff000  bootpack header 副本
0x40200000  kernel（页分配器从这里往后）
\end{txtcode}

打包由 \file{miniBoot/tools/mkbootpack.c} 完成：

\begin{shcode}
mkbootpack bootpack.bin kernel=minios-kernel.bin@0x40200000 \
    boot_xr829=.../boot_xr829.bin@0x40100000 \
    fw_xr829=.../fw_xr829.bin@0x40110000 \
    sdd_xr829=.../sdd_xr829.bin@0x40140000
\end{shcode}

顶层 \file{rpi\_bm/Makefile} 会自动处理：\code{firmware/xr829/} 下 boot/fw/sdd 三个文件齐全时填入 \code{BOOTPACK\_EXTRA}，否则打印 \code{WiFi firmware not packed} 并只打包 kernel。安装时 Boot0 写 8 KiB、bootpack 写 128 KiB：

\begin{shcode}
make t113
make install T113_DISK=/dev/目标设备
\end{shcode}

\section{图解：数据流}

\subsection{发送：一个 UDP 包到空口}

\begin{txtcode}
用户态 tftp / ping
  | svc：T113_SYS_TFTP / T113_SYS_NETCTL
  v
udp_sendto(s, ip, port, data, len)            net/ipv4/udp.c，加 rtnl_lock
  | skb = alloc_skb(len)                       头部预留 NET_SKB_PAD = 128 B
  | skb_push(8)                                UDP 头
  v
ip_queue_xmit(skb, daddr, IPPROTO_UDP)        net/ipv4/ip_output.c
  | skb_push(20)                               IP 头 + 校验和
  | arp_resolve(dev, 下一跳, mac)              表里没有就发 ARP 请求并返回 -EBUSY
  | eth_header()：skb_push(14)                 以太网头
  v
dev_queue_xmit(skb)                           net/core/dev.c，检查 carrier
  | dev->netdev_ops->ndo_start_xmit
  v
ieee80211_subif_start_xmit(skb, wlan0)        net/mac80211/tx.c
  | 取出目的 MAC 和 ethertype，skb_pull(14)
  | ieee80211_build_data_hdr()
  |   +-- LLC/SNAP 8 B（AA AA 03 00 00 00 + ethertype）
  |   +-- 有 PTK 时头部预留 CCMP IV 8 B、尾部预留 MIC 8 B，置 Protected 位
  |   +-- 802.11 头 24 B：ToDS，A1 = BSSID，A2 = 本机，A3 = 目的
  v
ieee80211_tx_skb(local, skb, AC_BE, ...)
  | cookie = ++next_cookie，登记 tx_status[cookie % N] = PENDING
  | local->hw.ops->tx
  v
xradio_tx(hw, skb)                            xradio/txrx.c
  | AC → WSM 队列；数据帧用策略 1、速率表 0x00a4
  v
wsm_tx_frame(cookie, frame, len, queue, ...)    xradio/wsm.c
  | 拷进 tx_buf：24 B wsm_tx 头 + 802.11 帧        （之后 mac80211 释放 skb）
  v
xradio_bh_tx(tx_buf, len, cap)                xradio/bh.c
  | 等固件输入缓冲，写 len 和序号，对齐
  v
sdio_memcpy_toio(func, HIF_ADDR(buf_id_tx, IN_OUT_QUEUE), buf, wlen)
  | sdio_io_rw_extended()：CMD53，写，function 1，地址自增
  v
smhc_request() → smhc_pio()                   SMHC1：CPU 逐字写 FIFO（0x200），1-bit 12 MHz
  v
XR829 固件：用 TK 做 CCMP 加密、算 MIC 和 FCS，按速率表发送和重试
  |
  +-- TX confirm 0x0404 {packetID = cookie, status}
        → 下一次 xradio_bh_poll() 收到 → xradio_tx_confirm_cb()
        → ieee80211_tx_status(hw, cookie, status)     记到 tx_status 槽
\end{txtcode}

一路上 skb 只在两个地方被拷贝：用户数据进 skb 一次，skb 进 \code{tx\_buf} 一次（因为 CMD53 要求 4 字节对齐的连续缓冲，而且长度要按块对齐）。

\subsection{接收：从空口到 socket}

\begin{txtcode}[basicstyle=\ttfamily\scriptsize]
XR829 固件：收帧、查 BSSID 过滤、用 TK/GTK 解密（保留 IV 和 MIC）
  | RX indication 0x0804 放进输出队列，CONTROL.NEXT_LEN ≠ 0
  v
（谁在轮询）udp_recvfrom / icmp_ping / wpa 握手 / net_poll
  +-- __net_poll(dev, 10)                         net/core/dev.c
      +-- ndo_poll = ieee80211_poll()             net/mac80211/main.c
          +-- ops->poll = xradio_poll() → xradio_bh_poll(10)        【底半部】
          |   +-- rx_one()
          |       +-- CMD53 读 IN_OUT_QUEUE      SMHC1 PIO 读 FIFO
          |       +-- wsm_handle_rx(msg)
          |           +-- xradio_rx_cb(msg, len)                     xradio/txrx.c
          |               +-- status ≠ 0 → 丢弃
          |               +-- alloc_skb，拷贝 802.11 帧
          |               +-- AES：头部后移 8 B 盖住 IV，尾部去掉 8 B MIC
          |               +-- rx_status：channel、signal、DECRYPTED | IV_STRIPPED | MIC_STRIPPED
          |               +-- ieee80211_rx_irqsafe(hw, skb)          只入 local->rx_queue
          +-- ieee80211_rx_tasklet(local)                           【tasklet】
              +-- ieee80211_rx_one(skb)
                  +-- 管理帧
                  |   +-- 扫描中 → ieee80211_scan_rx() → cfg80211 BSS 表
                  |   +-- 已加入 → ieee80211_rx_mgmt()：auth/assoc 响应、deauth
                  +-- 数据帧 → ieee80211_rx_data()
                      +-- 只收 FromDS 且来自本 BSSID 的
                      +-- 加密帧必须已被硬件解密
                      +-- ieee80211_data_to_8023()    802.11 + LLC → 以太网帧
                      +-- EAPOL（0x888e）→ cfg80211_rx_control_port() → wpa_rx_eapol()
                      +-- 其他：端口已授权才收
                          +-- eth_type_trans() → netif_rx(skb)       入 backlog
      +-- net_rx_action()                                            【软中断】
          +-- deliver_skb()：按 skb->protocol 查 packet_type
              +-- ETH_P_ARP → arp_rcv()
              +-- ETH_P_IP  → ip_rcv()
                  +-- ICMP → icmp_rcv()
                  +-- UDP  → udp_rcv() → 按端口挂到 udp_sock.rx_queue
      +-- run_net_work()                                             【工作队列】
          +-- 例如 supplicant 的 key_work：M4 的 TX status 回来后装钥

udp_recvfrom() 看到 rx_queue 非空 → 出队拷给调用者
\end{txtcode}

三个层次对应 Linux 的三种上下文。没有中断时，它们由同一次 \code{\_\_net\_poll()}按顺序跑完：

\begin{table}[H]
\centering
\small
\begin{tabularx}{\linewidth}{l>{\raggedright\arraybackslash}X>{\raggedright\arraybackslash}X>{\raggedright\arraybackslash}X}
\toprule
层次 & Linux & miniOS 里的位置 & 允许做什么 \\
\midrule
底半部 & SDIO 中断 → xradio bh 线程 & \code{xradio\_bh\_poll()} & 读队列、分派；回调只能入队或记标志 \\
tasklet & mac80211 RX tasklet & \code{ieee80211\_rx\_tasklet()} & 处理 802.11 帧，可以发帧、下命令 \\
软中断 & NET\_RX softirq & \code{net\_rx\_action()} & 交给 IP/ARP，转到 socket \\
工作队列 & workqueue & \code{run\_net\_work()} & 需要等待的事，例如装钥 \\
\bottomrule
\end{tabularx}
\end{table}

\section{调用链：连接一个 WPA2 AP}

\begin{txtcode}[basicstyle=\ttfamily\scriptsize]
register_netdev(wlan0) → NETDEV_REGISTER → wpa_netdev_event()
（或用户命令 wifi → T113_SYS_NETCTL → NETCTL_WIFI）
+-- wpa_supplicant_connect(dev)                                   net/wpa/supplicant.c
    +-- rtnl_lock()
    +-- __wpa_connect(dev)
    |   +-- cfg80211_disconnect()；清 IP
    |   +-- load_config()               读 /boot/wifi.txt：ssid、psk、可选静态 IP
    |   +-- cfg80211_scan(dev, 信道 1–13，50–110 ms)
    |   |   +-- ieee80211_scan() → xradio_hw_scan()             xradio/scan.c
    |   |       +-- wsm_scan() → START_SCAN 0x0007
    |   |           beacon / probe resp → RX → ieee80211_scan_rx() → BSS 表
    |   |           SCAN_COMPLETE 0x0806 → ieee80211_scan_completed()
    |   +-- select_bss()                按 wifi.txt 的 SSID 查 BSS 表，没配置时试内置默认 SSID
    |   +-- build_rsn()                 CCMP 单播，CCMP/TKIP 组播，PSK
    |   +-- wpa_pbkdf2(psk, ssid) → PMK；eapol_init()
    |   +-- cfg80211_connect(dev, {bss, RSN IE})                   net/wireless/sme.c
    |   |   +-- ieee80211_mgd_auth()                              net/mac80211/mlme.c
    |   |   |   +-- bss_info_changed(BSSID) → xradio_do_join()
    |   |   |   |   +-- 预加入 MIB：电源模式、BA 策略、PMF 策略
    |   |   |   |   +-- wsm_join()              JOIN 0x000b：信道、BSSID、SSID、基本速率
    |   |   |   |   +-- 过滤 MIB：RX 过滤、beacon 过滤表/开关、BSSID 过滤
    |   |   |   +-- 发 Open System auth（管理帧，最低速率），等 auth 响应
    |   |   +-- ieee80211_mgd_assoc()
    |   |       +-- 发 assoc request（带 RSN IE），等 assoc 响应
    |   |       +-- bss_info_changed(ASSOC) → xradio_set_assoc()
    |   |           +-- 关联模式、keep-alive、BSS 参数（AID）、beacon 唤醒周期、PM、slot
    |   +-- 等四次握手（__net_poll 循环，最长 HANDSHAKE_TIMEOUT）
    |       M1 ← EAPOL 经 control port → eapol.c 生成 SNonce，算 PTK
    |       M2 → cfg80211_tx_control_port() → ieee80211_tx_control_port()
    |       M3 ← 校验 MIC，解出 GTK
    |       M4 → 发送；TX status 回来后 net_schedule_work(key_work)
    |       key_work → install_keys()
    |           +-- cfg80211_add_key() → ieee80211_cfg_add_key() → xradio_set_key()
    |               +-- wsm_add_key_pairwise()   ADD_KEY 0x000c，槽 0，AES 成对钥
    |               +-- wsm_add_key_group()      槽 1，GTK + RSC
    +-- cfg80211_port_authorized()     打开数据端口，carrier on，打印 "wlan0: link up"
    +-- rtnl_unlock()
    +-- wifi.txt 有静态 IP → inet_set_addr()；否则 ipconfig_dhcp()
\end{txtcode}

\section{T113 实现与 Linux xradio 驱动的区别}

\begin{table}[H]
\centering
\small
\begin{tabularx}{\linewidth}{l>{\raggedright\arraybackslash}X>{\raggedright\arraybackslash}X}
\toprule
方面 & Linux（Allwinner xradio + mac80211） & T113 miniOS \\
\midrule
事件来源 & SDIO 卡中断（DAT1）或 \code{WL\_WAKE\_AP} GPIO 中断，唤醒 bh 线程 & 没有中断。bh 在 \code{ieee80211\_poll()}、等 confirm、等缓冲时被调用 \\
总线 & 4-bit，可到 50 MHz 高速模式 & 1-bit，12 MHz（HOSC 24 MHz 固定 /2，是 SMHC 源时钟上限） \\
数据搬运 & SMHC 内部 DMA（IDMAC） & CPU 读写 FIFO \\
固件来源 & \code{request\_firmware()}，从 \file{/lib/firmware} & \code{t113\_vfs\_open("/boot/*.bin")}，从 SD 卡 FAT32 分区；失败时用 SPL 预加载的 bootfs \\
固件下载时机 & probe 很快返回，下载放在工作队列 & 在 \code{xradio\_sdio\_probe()} 里同步完成，此时持有驱动核心锁 \\
bh & 内核线程，等待队列 + 原子计数 & 普通函数 \code{xradio\_bh\_poll()}，单线程，\code{in\_rx} 防重入 \\
命令等待 & \code{wait\_event\_timeout()} 睡眠 & 循环 \code{xradio\_bh\_poll(5)} 忙等 \\
TX 队列 & 每个 AC 一个队列，按 link 流控、可重排 & 不排队，直接写；只用 \code{hw\_bufs\_used} 节流 \\
虚拟接口 & 最多两个（STA + AP/P2P） & 只有 STA，接口号 0 \\
省电 & PS-Poll、动态休眠、SDIO 睡眠唤醒 & 一直保持 WUP，PM 关闭 \\
加密 & 硬件加解密 & 同样交给固件；驱动在 RX 时剥 IV/MIC，TX 时 mac80211 预留 IV 和 MIC 的位置 \\
802.11n & HT、A-MPDU、40 MHz & 只做 11b/g 关联，BA 策略写 0 \\
\bottomrule
\end{tabularx}
\end{table}

\section{预期串口输出}

（地址、CRC、HIF 寄存器值和 WSM 数据以实测为准，下面是格式示例）

\begin{txtcode}
bootpack: kernel @0x40200000 bytes=0x????????
bootpack: boot_xr829 @0x40100000 bytes=0x00000888
bootpack: fw_xr829 @0x40110000 bytes=0x0002aa20
bootpack: sdd_xr829 @0x40140000 bytes=0x000002e8
sdc1: SMHC host ready (card scan deferred to rescan)
mmc1: SDIO card rca 0x0001 OCR 0xffff0000 funcs 1 vendor 0x0a9e device 0x2282
mmc1: func 1 class 0x0000 vendor 0x0a9e device 0x2282
mmc1:0001:1: XRadio WiFi 0x0a9e:0x2282 enabled (IOEx 0x0002 IORx 0x0002 blksz 0x0200)
mmc1:0001:1: HIF_CONFIG 0x???????? (stable) HIF_CONTROL 0x????????
mmc1:0001:1: CMD53 write/read-back OK (HIF_TSET_GEN, 3 patterns)
mmc1:0001:1: SDIO clock 12000 kHz
mmc1:0001:1: SDD ref clock 26000 kHz, DPLL 0x0ec4f121
mmc1:0001:1: device awake
mmc1:0001:1: bootloader loaded (2184 bytes)
mmc1:0001:1: bootloader answered I_AM_HERE
mmc1:0001:1: firmware 32 KiB
mmc1:0001:1: firmware 64 KiB
mmc1:0001:1: firmware 96 KiB
mmc1:0001:1: firmware 128 KiB
mmc1:0001:1: firmware 160 KiB
mmc1:0001:1: firmware download complete (174624 bytes)
mmc1:0001:1: firmware running
wsm: firmware startup: ? x ???? B input buffers, hw ?.?, type ?, api 0x????????
wsm: firmware label "..."
mmc1:0001:1: WSM command layer ready
mmc1:0001:1: configure firmware, SDD bytes=0x000002e8
mmc1:0001:1: configured MAC 02:0a:9e:22:82:01 bands=0x???????? rates=0x????????
\end{txtcode}

\section{排障提示}

启动阶段：

\begin{itemize}
  \item \code{mmc1: no SDIO card (CMD5 timeout)}——先查 WL\_REG\_ON 电平和 32k fanout 是否真的在输出（\code{smhc\_dump()} 会打印 FANOUT\_GATE/CFG 和 PG 组引脚配置）；
  \item \code{card OCR has no 2.7-3.6V window}——R4 回来了但 OCR 位图异常，多半是总线信号质量（CLK 上没串 33R 就容易出这个）；
  \item \code{card never became ready (CMD5)}——模组在初始化，可适当加大 \code{power\_on\_delay\_ms}；
  \item \code{no sdio driver matched (N registered)}——vendor/device 不在 \code{xradio\_sdio\_ids[]} 里，看上一行打印的实测 ID；
  \item \code{cannot load /boot/...\_xr829.bin}——\file{/boot} 下缺固件：把三个 \code{*\_xr829.bin} 拷到 SD 卡的 FAT32 分区（\code{make install-firmware}），启动后用 \code{ls /boot} 确认；如果串口打印了 \code{fat32: mount failed}，再检查 bootpack 是否打包了固件。
\end{itemize}

CMD53 与固件下载：

\begin{table}[H]
\centering
\begin{tabular}{p{7.0cm}p{7.6cm}}
\toprule
打印 & 含义 \\
\midrule
\code{CMD53 not working; firmware download impossible} & HIF 寄存器读/写失败；查 pinmux、PIO FIFO 和总线时钟 \\
\code{CMD53 fails at the fast clock, back to the init clock} & 12 MHz 下总线不稳，已回退 400 kHz，下载变慢但仍可完成 \\
\code{not an XRadio HIF (CONFIG)} & CONFIG 的 hw-type 位不是 1：芯片/固件不匹配或访问错窗口 \\
\code{DPLL register did not stick} & TSET\_GEN 写入没生效，检查间接/直连访问模式 \\
\code{device did not wake up, CONTROL} & WUP 后 RDY 未置位：DPLL 或 32k 时钟问题 \\
\code{DPLL access timeout, COMMON\_CFG} & WLAN PWRCTRL 仲裁一直没拿到（BT 占用或时钟未起） \\
\code{bootloader not ready, IMAGE\_SIZE reads} & boot\_xr829 没跑起来：AHB 载入地址/CPU 释放/镜像损坏 \\
\code{bootloader reported error} & 设备侧 downloader 拒收，看 STATUS 值 \\
\code{timeout waiting for download FIFO} & GET 长时间不推进，设备侧卡住 \\
\code{firmware download failed, status ... trace pc ...} & 下载失败，TRACE\_PC 是设备侧出错位置 \\
\code{wsm: no startup indication from firmware} & 固件下载成功但没发 STARTUP indication，通常是 fw/sdd 版本不匹配 \\
\bottomrule
\end{tabular}
\end{table}

\section{当前边界与下一步}

硬件与固件链路（本章主线）已实现：

\begin{itemize}
  \item SDIO 枚举、CCCR/FBR/CIS、CMD52；
  \item CMD53 字节/块模式 + SMHC PIO 数据阶段；
  \item HIF 寄存器、SRAM/AHB 间接寻址、prefetch；
  \item SDD 参考时钟解析、DPLL 配置、设备唤醒；
  \item bootloader 载入、1 KiB 分块固件下载、下载 FIFO 协议；
  \item queue 模式切换、HIF 队列收发（bh.c）与 WSM 协议（命令/confirm/indication、MIB）。
\end{itemize}

协议与数据面（见第~\ref{chap:netstack} 章）也已实现：硬件扫描、WSM JOIN、认证/关联、WPA2-PSK 四次握手与装钥、802.11 $\leftrightarrow$ 以太网数据帧转换、ARP/IPv4/ICMP/UDP/DHCP、tftp 客户端与用户命令。

未实现（按优先级）：

\begin{enumerate}
  \item SMHC/SDIO 中断（现在整条链路轮询，吞吐受限）；
  \item 可选：SMHC DMA、high-speed 时序、蓝牙固件（fw\_xr829\_bt.bin）与 ETF 测试镜像。
\end{enumerate}

其他刻意简化：卡/函数结构用静态池（\code{card\_pool[3]}、\code{func\_pool[21]}，外加 1 KiB 的静态下载缓冲）而不是分配器；所有等待用 \code{mdelay}/\code{udelay} 忙等；固件下载在 probe 内同步完成，尚未搬到工作队列。
